% Options for packages loaded elsewhere
% Options for packages loaded elsewhere
\PassOptionsToPackage{unicode}{hyperref}
\PassOptionsToPackage{hyphens}{url}
\PassOptionsToPackage{dvipsnames,svgnames,x11names}{xcolor}
%
\documentclass[
  11pt,
  letterpaper,
]{article}
\usepackage{xcolor}
\usepackage[margin=1in]{geometry}
\usepackage{amsmath,amssymb}
\setcounter{secnumdepth}{5}
\usepackage{iftex}
\ifPDFTeX
  \usepackage[T1]{fontenc}
  \usepackage[utf8]{inputenc}
  \usepackage{textcomp} % provide euro and other symbols
\else % if luatex or xetex
  \usepackage{unicode-math} % this also loads fontspec
  \defaultfontfeatures{Scale=MatchLowercase}
  \defaultfontfeatures[\rmfamily]{Ligatures=TeX,Scale=1}
\fi
\usepackage{lmodern}
\ifPDFTeX\else
  % xetex/luatex font selection
\fi
% Use upquote if available, for straight quotes in verbatim environments
\IfFileExists{upquote.sty}{\usepackage{upquote}}{}
\IfFileExists{microtype.sty}{% use microtype if available
  \usepackage[]{microtype}
  \UseMicrotypeSet[protrusion]{basicmath} % disable protrusion for tt fonts
}{}
\makeatletter
\@ifundefined{KOMAClassName}{% if non-KOMA class
  \IfFileExists{parskip.sty}{%
    \usepackage{parskip}
  }{% else
    \setlength{\parindent}{0pt}
    \setlength{\parskip}{6pt plus 2pt minus 1pt}}
}{% if KOMA class
  \KOMAoptions{parskip=half}}
\makeatother
% Make \paragraph and \subparagraph free-standing
\makeatletter
\ifx\paragraph\undefined\else
  \let\oldparagraph\paragraph
  \renewcommand{\paragraph}{
    \@ifstar
      \xxxParagraphStar
      \xxxParagraphNoStar
  }
  \newcommand{\xxxParagraphStar}[1]{\oldparagraph*{#1}\mbox{}}
  \newcommand{\xxxParagraphNoStar}[1]{\oldparagraph{#1}\mbox{}}
\fi
\ifx\subparagraph\undefined\else
  \let\oldsubparagraph\subparagraph
  \renewcommand{\subparagraph}{
    \@ifstar
      \xxxSubParagraphStar
      \xxxSubParagraphNoStar
  }
  \newcommand{\xxxSubParagraphStar}[1]{\oldsubparagraph*{#1}\mbox{}}
  \newcommand{\xxxSubParagraphNoStar}[1]{\oldsubparagraph{#1}\mbox{}}
\fi
\makeatother


\usepackage{longtable,booktabs,array}
\newcounter{none} % for unnumbered tables
\usepackage{calc} % for calculating minipage widths
% Correct order of tables after \paragraph or \subparagraph
\usepackage{etoolbox}
\makeatletter
\patchcmd\longtable{\par}{\if@noskipsec\mbox{}\fi\par}{}{}
\makeatother
% Allow footnotes in longtable head/foot
\IfFileExists{footnotehyper.sty}{\usepackage{footnotehyper}}{\usepackage{footnote}}
\makesavenoteenv{longtable}
\usepackage{graphicx}
\makeatletter
\newsavebox\pandoc@box
\newcommand*\pandocbounded[1]{% scales image to fit in text height/width
  \sbox\pandoc@box{#1}%
  \Gscale@div\@tempa{\textheight}{\dimexpr\ht\pandoc@box+\dp\pandoc@box\relax}%
  \Gscale@div\@tempb{\linewidth}{\wd\pandoc@box}%
  \ifdim\@tempb\p@<\@tempa\p@\let\@tempa\@tempb\fi% select the smaller of both
  \ifdim\@tempa\p@<\p@\scalebox{\@tempa}{\usebox\pandoc@box}%
  \else\usebox{\pandoc@box}%
  \fi%
}
% Set default figure placement to htbp
\def\fps@figure{htbp}
\makeatother


% definitions for citeproc citations
\NewDocumentCommand\citeproctext{}{}
\NewDocumentCommand\citeproc{mm}{%
  \begingroup\def\citeproctext{#2}\cite{#1}\endgroup}
\makeatletter
 % allow citations to break across lines
 \let\@cite@ofmt\@firstofone
 % avoid brackets around text for \cite:
 \def\@biblabel#1{}
 \def\@cite#1#2{{#1\if@tempswa , #2\fi}}
\makeatother
\newlength{\cslhangindent}
\setlength{\cslhangindent}{1.5em}
\newlength{\csllabelwidth}
\setlength{\csllabelwidth}{3em}
\newenvironment{CSLReferences}[2] % #1 hanging-indent, #2 entry-spacing
 {\begin{list}{}{%
  \setlength{\itemindent}{0pt}
  \setlength{\leftmargin}{0pt}
  \setlength{\parsep}{0pt}
  % turn on hanging indent if param 1 is 1
  \ifodd #1
   \setlength{\leftmargin}{\cslhangindent}
   \setlength{\itemindent}{-1\cslhangindent}
  \fi
  % set entry spacing
  \setlength{\itemsep}{#2\baselineskip}}}
 {\end{list}}
\usepackage{calc}
\newcommand{\CSLBlock}[1]{\hfill\break\parbox[t]{\linewidth}{\strut\ignorespaces#1\strut}}
\newcommand{\CSLLeftMargin}[1]{\parbox[t]{\csllabelwidth}{\strut#1\strut}}
\newcommand{\CSLRightInline}[1]{\parbox[t]{\linewidth - \csllabelwidth}{\strut#1\strut}}
\newcommand{\CSLIndent}[1]{\hspace{\cslhangindent}#1}



\setlength{\emergencystretch}{3em} % prevent overfull lines

\providecommand{\tightlist}{%
  \setlength{\itemsep}{0pt}\setlength{\parskip}{0pt}}



 


\usepackage{booktabs}
\usepackage{microtype}
\makeatletter
\@ifpackageloaded{caption}{}{\usepackage{caption}}
\AtBeginDocument{%
\ifdefined\contentsname
  \renewcommand*\contentsname{Table of contents}
\else
  \newcommand\contentsname{Table of contents}
\fi
\ifdefined\listfigurename
  \renewcommand*\listfigurename{List of Figures}
\else
  \newcommand\listfigurename{List of Figures}
\fi
\ifdefined\listtablename
  \renewcommand*\listtablename{List of Tables}
\else
  \newcommand\listtablename{List of Tables}
\fi
\ifdefined\figurename
  \renewcommand*\figurename{Figure}
\else
  \newcommand\figurename{Figure}
\fi
\ifdefined\tablename
  \renewcommand*\tablename{Table}
\else
  \newcommand\tablename{Table}
\fi
}
\@ifpackageloaded{float}{}{\usepackage{float}}
\floatstyle{ruled}
\@ifundefined{c@chapter}{\newfloat{codelisting}{h}{lop}}{\newfloat{codelisting}{h}{lop}[chapter]}
\floatname{codelisting}{Listing}
\newcommand*\listoflistings{\listof{codelisting}{List of Listings}}
\makeatother
\makeatletter
\makeatother
\makeatletter
\@ifpackageloaded{caption}{}{\usepackage{caption}}
\@ifpackageloaded{subcaption}{}{\usepackage{subcaption}}
\makeatother
\usepackage{bookmark}
\IfFileExists{xurl.sty}{\usepackage{xurl}}{} % add URL line breaks if available
\urlstyle{same}
\hypersetup{
  pdftitle={Deliberation by Design: A Multi-Agent Large-Language-Model System for Constructively Aligned{,} Provenance-Tracked Examination Authoring in the Sciences},
  pdfauthor={Scott Mutchler},
  pdfkeywords={automatic item generation, large language
models, multi-agent systems, educational assessment, constructive
alignment, retrieval-augmented generation, provenance, accessibility},
  colorlinks=true,
  linkcolor={blue},
  filecolor={Maroon},
  citecolor={Blue},
  urlcolor={Blue},
  pdfcreator={LaTeX via pandoc}}


\title{Deliberation by Design: A Multi-Agent Large-Language-Model System
for Constructively Aligned, Provenance-Tracked Examination Authoring in
the Sciences}
\author{Scott Mutchler}
\date{2026-06-06}
\begin{document}
\maketitle
\begin{abstract}
Authoring a high-quality summative examination requires the simultaneous
satisfaction of constraints that are individually well understood but
jointly difficult to hold in mind: content coverage in the proportions a
course promises, cognitive demand matched to stated learning outcomes,
freedom from item-writing flaws, accessibility for students with diverse
needs, and defensible measurement properties. A single author working
under deadline tends to collapse these constraints into whichever is
easiest to satisfy, producing exams that drift toward recall, harbour
cueing flaws, and embed construct-irrelevant difficulty. We present an
agentic system that reframes exam construction as \emph{structured
disagreement} rather than single-pass generation. A team of specialised
large-language-model (LLM) agents --- a subject-matter expert, a
blueprint architect, and a panel of adversarial critics for alignment,
item-writing mechanics, accessibility, psychometrics, and test-wiseness
--- propose, critique, and refine items across bounded epochs,
coordinated by a \emph{deterministic} (non-LLM) moderator that applies a
faculty-specified trade-off policy and records every proposal,
objection, and resolution as an immutable provenance event. Item
proposals are grounded in the instructor's own uploaded materials by
mandatory retrieval citations and an independent grounding check, so the
exam is built from course content rather than from the model's
parametric memory. We describe the design, its grounding in established
assessment theory (backward design, Bloom's revised taxonomy,
constructive alignment, and evidence-centred design), and an end-to-end
reference implementation that produces a complete, render-ready exam
bundle --- exam, answer key, rubrics, accessibility variants, a coverage
and difficulty report, and a full audit trail --- in roughly three
minutes of wall-clock time on a two-tier model routing (a strong hosted
model for generative roles and a small hosted model for verification
roles); the same pipeline can be run entirely on local open-weight
models at the cost of speed. We discuss the pedagogical, cost, and time
trade-offs of the approach, including its limitations (the absence of
empirical item-response data, the residual risk of fluent-but-wrong
content, and the gap between estimated and observed psychometrics), and
we describe our plan to release the system as an open-source project to
invite scrutiny, replication, and disciplinary refinement.
\end{abstract}


\section{Introduction}\label{introduction}

Summative assessment is among the highest-stakes design tasks an
instructor performs, and among the least resourced. A midterm or final
examination is the instrument through which a course converts months of
instruction into defensible inferences about what students know and can
do. Yet the typical author of that instrument is a single faculty
member, working under deadline, holding in mind at once a set of
constraints that the educational-measurement literature treats as
distinct specialities: the exam must cover the right topics in the right
proportions; it must target the cognitive levels the course outcomes
actually promise rather than the levels that are easiest to write items
for; it must be free of the grammatical cueing, implausible distractors,
and overlapping options that the item-writing literature has catalogued
for decades (Haladyna et al. 2002); it must remain accessible to
students with diverse needs and language backgrounds, rather than
penalising them through construct-irrelevant difficulty (Messick 1995);
and it must distinguish students who have understood the material from
those who have not, within a fixed time and point budget.

Each of these dimensions has a mature literature behind it. Backward
design (Wiggins and McTighe 2005) tells us to specify the evidence an
assessment must elicit \emph{before} writing any tasks. Bloom's revised
taxonomy (Anderson and Krathwohl 2001) gives a vocabulary for cognitive
demand. Constructive alignment (Biggs 1996) requires that each task
elicit evidence at the cognitive level its target outcome promises.
Evidence-centred assessment design (Mislevy et al. 2003) frames every
item as an argument linking a claim about the student to the evidence a
task elicits. The \emph{Standards for Educational and Psychological
Testing} (American Educational Research Association et al. 2014) codify
validity and fairness as professional obligations. Few practising
faculty, however, have the time to consult these literatures for each
assessment they write, and the predictable consequence is a
well-documented set of failure modes: silent drift toward recall and
procedural application, item-writing flaws that inflate scores without
measuring anything, and accommodations produced reactively as ad hoc
edits to a finished exam rather than designed in from the start.

The recent maturation of large language models (LLMs) ({Brown et al.}
2020) has prompted enthusiasm for automating assessment authoring, and a
substantial body of work on automatic question generation predates the
current generation of models (Kurdi et al. 2020). The naïve application
of a single, powerful model to ``write me an exam,'' however, reproduces
the single-author failure mode in a new guise. Asked to optimise
content, cognition, accessibility, and measurement simultaneously, a
single model tends to produce a smooth compromise that \emph{looks}
plausible on every dimension while obscuring where the trade-offs were
made --- an argument we make on design grounds rather than from a
controlled comparison. Worse, LLMs are prone to confident fabrication
(Ji et al. 2023) and to sycophantic agreement that simulates
deliberation without its substance ({Sharma et al.} 2023; {Perez et al.}
2023). A generator that hallucinates a plausible-but-wrong answer key,
or that agrees with every critique it receives, is not merely unhelpful;
it is dangerous in a context where the output drives grades and is
subject to appeal.

This paper presents a system designed around the opposite premise. We
treat exam construction not as a generation problem but as a
\emph{deliberation} problem, and we make the deliberation structural. A
small team of specialised LLM agents --- a subject-matter expert (SME),
a blueprint architect, and a panel of adversarial critics --- propose
and contest items across bounded epochs. A \textbf{deterministic}
moderator, not itself an LLM, collects objections, applies a trade-off
policy the faculty member specifies in advance, and records every move
as an immutable provenance event. Item content is \textbf{grounded} in
the instructor's own uploaded materials through mandatory retrieval
citations and an independent grounding check, so the exam is assembled
from course content rather than from the model's parametric memory. The
instructor of record is retained as the authoritative decision-maker,
reviewing a fully rendered bundle rather than mediating the machinery.

Our contributions are: (1) a system design that operationalises four
established assessment frameworks as hard constraints in an agentic
pipeline, rather than as advisory prompts; (2) a deliberately
\emph{non-LLM} moderator and a mandatory-objection protocol that
together counter the multi-agent failure mode of
convergence-through-agreement; (3) a grounding regime that ties every
item to cited source passages and audits answer-key support; (4) a
full-provenance record engineered for defensibility against grade
appeals; and (5) a working, end-to-end reference implementation, which
we describe and which we are preparing to release as open source. We are
explicit throughout about what the system does \emph{not} yet do ---
most importantly, it produces \emph{estimated} rather than empirically
observed psychometrics --- and we frame the open-source release as the
mechanism by which those gaps are most credibly closed.

\section{Related Work}\label{related-work}

\textbf{Automatic item and question generation.} Automatic question
generation (AQG) has a long pre-LLM history, comprehensively surveyed by
Kurdi et al. (2020), who reviewed 93 studies and noted that the field
had paid little attention to controlling difficulty, enriching question
structure, or generating feedback, precisely the dimensions a
blueprint-driven, critic-mediated design foregrounds. Classical AQG
systems were largely template- or syntax-driven and produced factual
recall questions; they rarely engaged the cognitive-alignment or
fairness concerns that dominate professional assessment practice. The
arrival of instruction-tuned LLMs ({Brown et al.} 2020) dramatically
expanded the fluency and range of generated items, and recent work has
applied LLMs to question generation and to Bloom-level-targeted item
writing. Our work differs in treating generation as only the first move
in a multi-party deliberation, and in subordinating the generator to an
explicit blueprint and a panel of critics.

\textbf{Multi-agent LLM systems.} A growing literature shows that
orchestrating several LLM instances --- debating, critiquing, or playing
distinct roles --- improves factuality and reasoning relative to a
single pass. Du et al. (2024) demonstrate that multi-agent debate
reduces hallucination and improves reasoning; Park et al. (2023) show
that role-specialised agents can sustain coherent, differentiated
behaviour; and frameworks such as AutoGen (Wu et al. 2023) provide
general machinery for multi-agent conversation. Our design draws on this
line of work but departs from it in two deliberate ways. First, the
coordinating moderator is \emph{deterministic} rather than a model: the
trade-off resolution that decides which dimension wins a conflict is a
priority-ranked policy the faculty member sets, not an emergent property
of a conversation. We made this choice because the resolution of a
content-versus- accessibility conflict is a \emph{value judgment} that
must be auditable and must belong to the instructor, not to an opaque
negotiation. Second, we adopt an anti-sycophancy protocol (critics must
produce objections at a configurable minimum rate in early epochs, and
empty critique is rejected and re-prompted) to counter the
well-documented tendency of LLMs toward agreeable convergence ({Sharma
et al.} 2023; {Perez et al.} 2023).

\textbf{Grounding and retrieval.} Retrieval-augmented generation (Lewis
et al. 2020) conditions model output on retrieved documents to improve
factual accuracy and reduce hallucination. We apply this idea as a
\emph{gating constraint} rather than a quality nudge: an item proposal
that does not cite at least one retrieved chunk from the instructor's
materials is rejected before it can enter the item bank, and an
independent grounding verifier checks that the cited passages actually
support the claimed answer. Prior examinations, when supplied, are
indexed for \emph{style} only and are never returned as a content source
for new items.

\textbf{Assessment theory.} The system is an attempt to make four
established frameworks operative rather than aspirational. Backward
design (Wiggins and McTighe 2005) is enforced by forbidding item
generation before a blueprint is approved. Bloom's revised taxonomy
(Anderson and Krathwohl 2001) is the vertical axis of that blueprint.
Constructive alignment (Biggs 1996) is checked per item by a dedicated
agent. Evidence-centred design (Mislevy et al. 2003), together with
modern validity theory (Messick 1995; Kane 2013), motivates the
provenance record that turns each item into a documented argument.
Accessibility is treated as a first-class design constraint in the
spirit of Universal Design for Learning (Rose and Meyer 2002; CAST 2018)
rather than as post-hoc accommodation.

\section{System Design}\label{system-design}

\subsection{Overview and design
commitments}\label{overview-and-design-commitments}

The system is delivered as a web application but its substance is a team
of agents that propose, critique, and refine items under instructor
supervision. Five commitments shape the design and distinguish it from a
prompt-engineered single-model generator.

\begin{enumerate}
\def\labelenumi{\arabic{enumi}.}
\tightlist
\item
  \textbf{Multi-agent, not single-agent.} Each critic --- alignment,
  item-writing, accessibility, psychometrics, adversarial student --- is
  a separate agent with its own persona and output schema. The point is
  to make trade-offs \emph{visible}, not to optimise them away.
\item
  \textbf{The moderator is deterministic.} Disagreement is resolved by a
  priority-ranked policy, not by a model. This keeps value judgments
  with the instructor and keeps resolution auditable.
\item
  \textbf{Grounding is mandatory.} Every item proposal must cite at
  least one source chunk from uploaded materials, and a grounding
  verifier checks that the cited chunks support the answer. Ungrounded
  items are auto-rejected. Prior exams are indexed for style only.
\item
  \textbf{Full provenance.} Every proposal, edit, objection, and
  resolution writes a structured provenance event. This is what makes an
  exam defensible under appeal.
\item
  \textbf{No agent-orchestration frameworks.} All agents and
  orchestration logic are written directly against provider SDKs with an
  explicit, inspectable state machine, so that the control flow ---
  particularly the deterministic moderator --- is fully legible and
  testable.
\end{enumerate}

\subsection{Inputs}\label{inputs}

The instructor provides three things. \textbf{Materials} are course
readings, lecture notes, problem sets, and lab manuals (PDF and related
formats); optional prior exams are accepted for style reference only. A
\textbf{course specification} lists course learning outcomes (CLOs),
each tagged with a target Bloom level, a weighted topic list, and
free-text guiding principles (e.g., ``emphasise quantitative reasoning;
no rote definitions''). An \textbf{exam specification} fixes the exam
type, total points and time budget, item-type counts, target difficulty
distribution, and any required accessibility variants. A fourth input
deserves emphasis: the \textbf{trade-off policy}, a ranked priority over
\emph{content fidelity, cognitive alignment, accessibility,
discrimination,} and \emph{brevity}, together with a maximum epoch count
and a convergence rule. Requiring the instructor to rank these in
advance places the responsibility for value judgments where it belongs
and prevents the system from silently defaulting to whichever dimension
is easiest to quantify.

\subsection{Agent roster}\label{agent-roster}

Table 1 lists the agents. Each is defined by a system prompt (its
``constitution''), a set of allowed actions, and a forced-JSON output
schema validated against a Pydantic contract.

\begin{longtable}[]{@{}
  >{\raggedright\arraybackslash}p{(\linewidth - 4\tabcolsep) * \real{0.3333}}
  >{\raggedright\arraybackslash}p{(\linewidth - 4\tabcolsep) * \real{0.3333}}
  >{\raggedright\arraybackslash}p{(\linewidth - 4\tabcolsep) * \real{0.3333}}@{}}
\caption{The agent roster. The moderator is the only non-LLM
participant.}\label{tbl-roster}\tabularnewline
\toprule\noalign{}
\begin{minipage}[b]{\linewidth}\raggedright
Agent
\end{minipage} & \begin{minipage}[b]{\linewidth}\raggedright
Role
\end{minipage} & \begin{minipage}[b]{\linewidth}\raggedright
Tier
\end{minipage} \\
\midrule\noalign{}
\endfirsthead
\toprule\noalign{}
\begin{minipage}[b]{\linewidth}\raggedright
Agent
\end{minipage} & \begin{minipage}[b]{\linewidth}\raggedright
Role
\end{minipage} & \begin{minipage}[b]{\linewidth}\raggedright
Tier
\end{minipage} \\
\midrule\noalign{}
\endhead
\bottomrule\noalign{}
\endlastfoot
Subject-Matter Expert (SME) & Extracts themes; proposes grounded items;
defends or revises under critique & High \\
Blueprint Architect & Builds the topic × Bloom matrix governing coverage
& High \\
Learning-Outcomes Alignment (LOA) & Verifies each item targets the
cognitive level its CLO promises & Low \\
Item-Writing Specialist (IWS) & Catches cueing, implausible distractors,
and other mechanical flaws & Low \\
Accessibility Expert & Reviews readability and construct-irrelevant
difficulty; builds variants & Low \\
Psychometrician & Estimates difficulty and discrimination; audits the
distribution & Low \\
Adversarial Student & Attempts each item using only test-wiseness, no
materials & High \\
Grounding Verifier & Checks cited chunks actually support the claimed
answer & Low \\
Moderator & Deterministic orchestrator; applies the trade-off policy &
--- (not an LLM) \\
\end{longtable}

The ``tier'' column reflects a deliberate cost decision discussed in
Section~\ref{sec-cost}: generative, reasoning-heavy roles (SME,
architect, adversarial student) are routed to a strong hosted model,
while verification and rule-application roles (alignment, item-writing,
accessibility, psychometrics, grounding) are routed to a small,
inexpensive model.

\subsection{Orchestration: the epoch
loop}\label{orchestration-the-epoch-loop}

The pipeline runs in four phases, gated in the original design by three
faculty checkpoints (CP1: blueprint approval; CP2: item-bank review;
CP3: final review).

\begin{itemize}
\tightlist
\item
  \textbf{Phase 0 --- Intake.} Materials are ingested, chunked with a
  heading-aware chunker, embedded, and stored in a local vector index.
  The specifications and trade-off policy are parsed.
\item
  \textbf{Phase 1 --- Blueprinting.} The SME extracts and ranks
  candidate themes from the materials; the Blueprint Architect builds a
  topic × Bloom matrix with target item counts and points per cell. The
  blueprint is the central artifact, and it is frozen at CP1 before any
  item is written --- the structural enforcement of backward design.
\item
  \textbf{Phase 2 --- Item generation.} For each blueprint cell, the SME
  \emph{over-generates} candidate items (each carrying mandatory source
  citations), the IWS performs a mechanics cleanup pass, and the LOA
  verifies constructive alignment, dropping misaligned items. When the
  LOA rejects an item for a Bloom mismatch, the system first attempts an
  automatic realignment (re-mapping the item's Bloom level or CLO
  reference and re-verifying) before discarding it.
\item
  \textbf{Phase 3 --- Refinement epochs.} In each epoch, the critic
  panel reviews every surviving item \emph{in parallel}: the
  Accessibility Expert flags readability and construct-irrelevant
  difficulty, the Adversarial Student attempts each item using only
  test-wiseness, and the Psychometrician scores difficulty and flags
  distributional imbalance. The moderator collects the objections,
  routes each to the SME for an edit or a rebuttal, applies the
  trade-off policy when critics conflict, and re-verifies any edited
  item through IWS → LOA → grounding. The loop exits when no critical-
  or high-severity objections are raised for a full epoch, or when the
  maximum epoch count is reached.
\item
  \textbf{Phase 4 --- Finalisation.} The Accessibility Expert generates
  the required accommodation variants; the Psychometrician produces an
  exam-level report (Bloom distribution, estimated difficulty curve,
  CLO-coverage heatmap); and the system renders the export bundle.
\end{itemize}

\subsection{The deterministic
moderator}\label{the-deterministic-moderator}

The moderator does the coordinating, and it is the one component we
deliberately kept \emph{out} of the language-model layer. For each open
objection it applies a fixed procedure: a critical-severity objection
forces an SME edit or item rejection; a conflict between objections on
the same item is resolved by consulting the instructor's priority
ranking and selecting the higher-ranked dimension, with the rationale
logged; everything else is routed to the SME for an accept-edit, rebut,
or defer response. Rebut--reassert cycles are bounded before escalation.
Every resolution writes a provenance event. Because this logic is
deterministic, it is unit-testable, reproducible, and legible to a
faculty member who wants to know exactly why a given item survived or
died --- properties that an LLM-mediated negotiation cannot offer.

\subsection{Anti-sycophancy by
construction}\label{anti-sycophancy-by-construction}

Multi-agent LLM systems are prone to convergence through mutual
agreement, a failure mode that produces the appearance of deliberation
without its substance ({Sharma et al.} 2023; {Perez et al.} 2023). We
counter this structurally. Critic personas are written to reward
substantive disagreement; the moderator's mandatory-objection protocol
rejects an empty early-epoch critique and re-prompts; and the
Adversarial Student operates only on item text, with no access to
materials or answer key, so its successes are genuine test-wiseness
exploits rather than informed solving. We note only an informal
observation from the handful of reference runs we have inspected so far
--- rebuttals are frequent and substantive rather than perfunctory ---
and we flag it as anecdotal, not a measured rate; quantifying objection
and rebuttal behaviour across many runs is future work.

\subsection{Grounding and provenance}\label{grounding-and-provenance}

Every SME item proposal must cite at least one retrieved chunk;
proposals without citations are auto-rejected, and a separate Grounding
Verifier confirms that the cited chunks support the claimed answer,
routing failures back to the SME. This is retrieval-augmented generation
(Lewis et al. 2020) used as a \emph{gate} rather than a hint. The
provenance record --- which agent proposed each item, which objections
were raised, how each was resolved, and against which source passages
--- is the artifact that operationalises evidence-centred design
(Mislevy et al. 2003) and modern validity theory (Kane 2013): it turns
each item into a documented argument and gives an instructor a concrete,
auditable record to put before a grade-appeal committee.

\section{Reference Implementation}\label{reference-implementation}

We built a complete, end-to-end implementation to validate the design
rather than to ship a product. It comprises roughly thirty Pydantic data
contracts, one agent class per persona (each loading its constitution
from an editable Markdown file), a provider abstraction supporting a
hosted model, an OpenAI-compatible institutional endpoint, and a local
model server, a retrieval layer over a local vector store with a local
embedding model, an append-only event log with per-call input/output
sidecars, the deterministic moderator state machine, and an export
module that renders the final bundle through Quarto. A five-page
observation UI lets a faculty member launch a run, watch an interactive
timeline of every agent call, inspect outputs, edit personas, and read
the documentation.

Two implementation realities are worth reporting because they generalise
to anyone building agentic pipelines on current models. First,
\textbf{model output is not reliably schema-conformant.} Models
stringify nested JSON fields, emit invalid escape sequences when
embedding LaTeX backslashes, and omit required fields when terse. We
addressed this principally with a retry-on-validation loop --- when a
response fails its Pydantic contract, the validation error is fed back
to the model and the call is retried --- backed by defensive defaults
and per-cell/per-item safe wrappers so that one malformed response can
never collapse a parallel fan-out. Second, \textbf{concurrency must be
governed per backend.} A single-GPU local model serialises internally
and times out under naïve parallelism; a free institutional endpoint has
an hourly quota. Per-provider semaphores (and a raised token floor for
local ``thinking'' models that emit long hidden reasoning streams before
producing content) were necessary for the parallel critique phase to
complete. We report these not as incidental engineering but as evidence
that robustness scaffolding is a first-class concern in agentic
assessment tooling.

On a two-tier routing --- a strong hosted model for the high tier and a
small hosted model for the low tier --- a quiz-scale reference corpus
runs the full Phase 0--4 pipeline, including four refinement epochs, in
roughly three minutes of wall-clock time, and produces a complete
bundle: the exam, answer key, instructor notes, rubrics, an exam report,
and a machine-readable provenance file, in PDF, DOCX, LaTeX, and HTML,
together with an on-demand human-readable narrative of the run. We
emphasise that this is a \emph{functional} demonstration on a small
corpus, not a controlled study of exam quality; the distinction is
central to the limitations we discuss below.

\section{Discussion: Trade-offs of the
Approach}\label{discussion-trade-offs-of-the-approach}

We assess the approach along the dimensions that matter to a department
deciding whether to adopt it, and we try not to flinch at the
unfavourable ones.

\subsection{Pedagogical trade-offs}\label{pedagogical-trade-offs}

The strongest pedagogical argument for the system is that it makes the
\emph{deliberative structure} of good assessment design unavoidable.
Backward design, constructive alignment, Bloom-level targeting, and
item-writing discipline are not suggestions layered on top of a
generator; they are enforced by the control flow. A faculty member
cannot get items before approving a blueprint, cannot keep a misaligned
item that the LOA rejects without an explicit override, and cannot ship
an item whose answer is unsupported by cited materials. This is the
educational-measurement equivalent of code review, and it is aimed
squarely at the kinds of failure modes the literature warns about
(Haladyna et al. 2002; Messick 1995), though whether it reduces them in
practice is an empirical question we have not yet answered. Treating
accessibility as a first-class critic, rather than a post-hoc
accommodation, aligns the workflow with Universal Design for Learning
(Rose and Meyer 2002).

The counter-argument is that the system can \emph{displace} rather than
\emph{develop} faculty expertise. A tool that silently fixes cueing
flaws does not teach the instructor to avoid them next time, and an
instructor who defers to the machine's blueprint may exercise less of
the curricular judgment the blueprint is meant to externalise. We have
tried to mitigate this by making the system observational and the
provenance legible --- the faculty member sees \emph{why} each decision
was made --- but the risk of deskilling is real and is best evaluated
empirically in a pilot. There is also a genuine pedagogical hazard
specific to LLMs: a fluent but subtly wrong item or answer key (Ji et
al. 2023) is more dangerous than an obviously wrong one, because it
survives a casual read. Mandatory grounding and the adversarial critic
reduce this risk but do not eliminate it, which is why the instructor
checkpoint on the rendered bundle is not optional in the intended
workflow.

\subsection{Cost trade-offs}\label{sec-cost}

Cost is governed almost entirely by model routing. The two-tier scheme
is the key economic lever: a typical run issues on the order of a few
dozen high-tier (generative) calls and roughly one hundred low-tier
(verification) calls. Routing the high tier to a strong hosted model and
the low tier to a small or local model concentrates spend on the
generative roles whose output most shapes the exam, while keeping the
marginal cost of an exam low. We have not benchmarked per-exam cost
directly; an order-of-magnitude estimate from the call counts and
current hosted-model token prices is single-digit cents to a few dollars
per exam, dominated by the high-tier calls. Marginal cost approaches
zero when the low tier (and optionally the high tier) is served by a
local model on departmental hardware, at the cost of speed and,
plausibly, some output quality. The system is provider-agnostic for
exactly this reason, and the local-only configuration buys an additional
benefit beyond cost: full data locality, since no course materials leave
the institution. The principal \emph{non-monetary} cost is engineering
and maintenance: agentic pipelines on current models require the
robustness scaffolding described above, and that scaffolding must track
model and SDK changes.

\subsection{Time trade-offs}\label{time-trade-offs}

The time argument is the most concrete. The system reallocates faculty
time from the mechanical aspects of construction --- drafting items,
checking distractors, balancing coverage, building accommodation
variants, assembling an answer key and rubrics --- to the consequential
ones: specifying outcomes, ranking trade-offs, and reviewing a finished
draft. The marginal machine time to produce a full bundle is minutes.
The faculty time is front-loaded into the course and exam
specifications, which, for a recurring course with stable outcomes,
amortises across semesters: the specification is written once and
reused, and only the materials and exam parameters change. The
break-even is least favourable for a one-off exam in a course whose
outcomes have never been formalised, and most favourable for a recurring
exam in a course with a stable, documented outcome set --- precisely the
qualifying-exam and large-enrolment-core settings where defensibility
matters most.

\subsection{Why multiple agents,
empirically}\label{why-multiple-agents-empirically}

A reasonable skeptic asks whether the multi-agent structure earns its
complexity. Our position, consistent with the multi-agent-debate
literature (Du et al. 2024), is that separating the critics is what
makes the trade-offs \emph{visible and contestable}. When the
Accessibility Expert objects that a stem carries construct-irrelevant
difficulty and the SME rebuts that the vocabulary in question \emph{is}
the construct, the disagreement is recorded, the policy is consulted,
and the resolution is logged with rationale. A single agent asked to
balance the same concerns internally produces a compromise with no
visible seam and no audit trail. The cost is more calls and more
orchestration; the benefit is accountability, which is what a graded,
appealable context demands.

\section{Limitations}\label{limitations}

The boundaries of the present work are worth stating plainly.

\textbf{Estimated, not observed, psychometrics.} The Psychometrician's
difficulty and discrimination figures are \emph{model estimates}, not
statistics derived from student responses. Real item parameters are
knowable only after administration. The system's psychometric report
should be read as a design-time sanity check on the \emph{intended}
distribution, not as a substitute for post-hoc item analysis. Closing
this gap requires a feedback loop in which administered-item statistics
are returned to the system --- work we have not yet done.

\textbf{Residual fabrication risk.} Grounding and adversarial critique
reduce but do not eliminate the chance of a fluent, plausible, and wrong
item or answer key (Ji et al. 2023). The mandatory instructor review of
the rendered bundle is the backstop, and the workflow depends on it: the
system is decision-support, not an autonomous authority.

\textbf{No controlled evaluation yet.} The implementation is verified
end-to-end on a small reference corpus. We have \emph{not} run a
controlled study comparing system-assisted exams to instructor practice,
nor measured inter-rater agreement between the agent critics and human
assessment experts, nor validated the constructive-alignment judgments
against expert raters. These are the studies that would establish the
system's value, and they are the central items of future work below.
Claims of quality improvement in this paper are therefore
\emph{architectural and theoretical}, not empirical.

\textbf{Scope.} The approach suits exams whose value lies in structured
coverage of stable outcomes --- recurring undergraduate exams, graduate
qualifying exams. It is ill-suited to assessments whose pedagogical
value is unpredictability (oral exams, open-ended research proposals),
and we make no claim for those settings.

\textbf{Generalisation of the corpus.} The reference runs use a small
science corpus. The theme-extraction batching path that would be
exercised by a full textbook is implemented but not yet stress-tested at
that scale, and discipline-specific conventions beyond the sciences
(e.g., the assessment norms of the humanities) are untested.

\textbf{Human-in-the-loop posture.} The original design specified
faculty checkpoints at CP1, CP2, and CP3. The reference implementation
auto-advances these and places the human review at the rendered-bundle
stage instead. This is a defensible UX choice for an observation-first
tool, but it weakens the in-loop control the design envisioned; making
the mid-run checkpoints optionally blocking is straightforward and is on
the roadmap.

\section{Future Work}\label{future-work}

Several directions follow directly from the limitations.

\begin{enumerate}
\def\labelenumi{\arabic{enumi}.}
\tightlist
\item
  \textbf{Empirical validation.} The priority is a pilot: a small number
  of faculty, one course and exam each, with a structured debrief
  comparing the system-assisted exam to prior practice on time, item
  quality, and faculty experience. In parallel, an expert-rater study
  would measure agreement between the agent critics (especially the
  LOA's alignment judgments) and human assessment specialists.
\item
  \textbf{Closing the psychometric loop.} Returning administered-item
  statistics --- difficulty, discrimination, distractor analysis --- to
  the system would let it calibrate its estimates and would convert the
  design-time report into a genuine measurement instrument over repeated
  administrations.
\item
  \textbf{Programmatic anti-sycophancy enforcement and policy-driven
  conflict resolution.} The mandatory-objection rate and the
  priority-rank consultation for directly conflicting critic demands are
  specified in the design and only partially enforced in code;
  completing them tightens the deliberation guarantees.
\item
  \textbf{Standards interoperability and persistence.} Export to QTI /
  Common Cartridge for direct LMS import, and a durable persistence
  layer for cross-session resume and longitudinal item banking, would
  move the system from a single-session tool toward a departmental
  asset.
\item
  \textbf{Broader disciplinary and accessibility coverage.}
  Stress-testing theme extraction on full textbooks, validating
  non-science disciplines, and expanding the accessibility variants
  beyond the current set are all tractable extensions.
\end{enumerate}

\section{Open-Source Release and
Collaboration}\label{open-source-release-and-collaboration}

We intend to release the system as an open-source project, and the
decision is methodological rather than merely generous. Three of the
system's commitments --- auditability, a deterministic and inspectable
moderator, and reproducibility --- are credible only if the code is open
to inspection. An assessment tool that asks a department to trust its
alignment judgments, while hiding how those judgments are made,
contradicts its own provenance argument. The deliberate avoidance of
heavyweight agent-orchestration frameworks in favour of a plain, legible
state machine was made partly with this in mind: a reviewer can read the
moderator and verify that the trade-off resolution does what the policy
says, without untangling a framework's control flow.

The release is also the most credible path to the empirical validation
the limitations demand. The studies that would establish the system's
worth --- pilots across courses and disciplines, expert-rater agreement,
psychometric calibration against administered data --- are larger than
any single group can run, and they depend on instrument transparency to
be comparable across sites. An open project with editable agent personas
(each is a Markdown file), provider-agnostic model routing, and a
documented data-contract layer invites exactly the multi-site,
multi-discipline contribution that turns a reference implementation into
a validated method. We plan to publish the repository with the reference
corpus, the agent constitutions, the deterministic moderator, and the
rendering pipeline; to document the contribution surfaces (new critic
agents, new export targets, new disciplinary persona sets); and to
license it for unrestricted academic use. We invite the assessment and
educational-technology communities to test the alignment judgments
against their own expertise, to contribute critics for concerns we have
not modelled, and to report where the system fails --- which, for a tool
whose entire premise is structured disagreement, is the most useful
contribution of all.

\section{Conclusion}\label{conclusion}

We have described a multi-agent LLM system that treats examination
authoring as structured deliberation rather than single-pass generation.
Four assessment frameworks --- backward design, Bloom's revised
taxonomy, constructive alignment, and evidence-centred design --- enter
the system as hard constraints rather than advice. A deterministic,
auditable moderator coordinates the critic agents under a trade-off
policy the instructor sets. Every item is grounded in cited course
materials and carries a full provenance trail. Taken together, these
make the trade-offs of good assessment design visible, contestable, and
defensible. A working reference implementation produces a complete,
render-ready exam bundle in minutes at low marginal cost, with the
instructor retained as the authoritative reviewer. The approach's
pedagogical, cost, and time trade-offs are favourable for recurring,
defensibility-sensitive assessments and unfavourable for one-off or
open-ended ones. Its central limitation --- estimated rather than
observed psychometrics, and the absence of controlled evaluation --- is
real, and we have argued that the most credible route to closing it is
an open-source release that invites the multi-site, multi-discipline
scrutiny on which assessment validity ultimately depends.

\section*{References}\label{references}
\addcontentsline{toc}{section}{References}

\protect\phantomsection\label{refs}
\begin{CSLReferences}{1}{1}
\bibitem[\citeproctext]{ref-aera2014standards}
American Educational Research Association, American Psychological
Association, and National Council on Measurement in Education. 2014.
\emph{Standards for Educational and Psychological Testing}. American
Educational Research Association.

\bibitem[\citeproctext]{ref-anderson2001taxonomy}
Anderson, Lorin W., and David R. Krathwohl. 2001. \emph{A Taxonomy for
Learning, Teaching, and Assessing: A Revision of Bloom's Taxonomy of
Educational Objectives}. Longman.

\bibitem[\citeproctext]{ref-biggs1996constructive}
Biggs, John. 1996. {``Enhancing Teaching Through Constructive
Alignment.''} \emph{Higher Education} 32 (3): 347--64.
\url{https://doi.org/10.1007/BF00138871}.

\bibitem[\citeproctext]{ref-brown2020gpt3}
{Brown, Tom B., Benjamin Mann, Nick Ryder, et al.} 2020. {``Language
Models Are Few-Shot Learners.''} \emph{Advances in Neural Information
Processing Systems (NeurIPS)} 33: 1877--901.

\bibitem[\citeproctext]{ref-cast2018udl}
CAST. 2018. \emph{Universal Design for Learning Guidelines Version 2.2}.
\href{https://udlguidelines.cast.org}{Https://udlguidelines.cast.org}.

\bibitem[\citeproctext]{ref-du2024multiagent}
Du, Yilun, Shuang Li, Antonio Torralba, Joshua B. Tenenbaum, and Igor
Mordatch. 2024. {``Improving Factuality and Reasoning in Language Models
Through Multiagent Debate.''} \emph{Proceedings of the 41st
International Conference on Machine Learning (ICML)}.

\bibitem[\citeproctext]{ref-haladyna2002review}
Haladyna, Thomas M., Steven M. Downing, and Michael C. Rodriguez. 2002.
{``A Review of Multiple-Choice Item-Writing Guidelines for Classroom
Assessment.''} \emph{Applied Measurement in Education} 15 (3): 309--34.
\url{https://doi.org/10.1207/S15324818AME1503_5}.

\bibitem[\citeproctext]{ref-ji2023survey}
Ji, Ziwei, Nayeon Lee, Rita Frieske, et al. 2023. {``Survey of
Hallucination in Natural Language Generation.''} \emph{ACM Computing
Surveys} 55 (12): 1--38. \url{https://doi.org/10.1145/3571730}.

\bibitem[\citeproctext]{ref-kane2013validating}
Kane, Michael T. 2013. {``Validating the Interpretations and Uses of
Test Scores.''} \emph{Journal of Educational Measurement} 50 (1): 1--73.
\url{https://doi.org/10.1111/jedm.12000}.

\bibitem[\citeproctext]{ref-kurdi2020systematic}
Kurdi, Ghader, Jared Leo, Bijan Parsia, Uli Sattler, and Salam Al-Emari.
2020. {``A Systematic Review of Automatic Question Generation for
Educational Purposes.''} \emph{International Journal of Artificial
Intelligence in Education} 30 (1): 121--204.
\url{https://doi.org/10.1007/s40593-019-00186-y}.

\bibitem[\citeproctext]{ref-lewis2020rag}
Lewis, Patrick, Ethan Perez, Aleksandra Piktus, et al. 2020.
{``Retrieval-Augmented Generation for Knowledge-Intensive {NLP}
Tasks.''} \emph{Advances in Neural Information Processing Systems
(NeurIPS)} 33: 9459--74.

\bibitem[\citeproctext]{ref-messick1995validity}
Messick, Samuel. 1995. {``Validity of Psychological Assessment:
Validation of Inferences from Persons' Responses and Performances as
Scientific Inquiry into Score Meaning.''} \emph{American Psychologist}
50 (9): 741--49. \url{https://doi.org/10.1037/0003-066X.50.9.741}.

\bibitem[\citeproctext]{ref-mislevy2003structure}
Mislevy, Robert J., Linda S. Steinberg, and Russell G. Almond. 2003.
{``Focus Article: On the Structure of Educational Assessments.''}
\emph{Measurement: Interdisciplinary Research and Perspectives} 1 (1):
3--62. \url{https://doi.org/10.1207/S15366359MEA0101_02}.

\bibitem[\citeproctext]{ref-park2023generative}
Park, Joon Sung, Joseph C. O'Brien, Carrie J. Cai, Meredith Ringel
Morris, Percy Liang, and Michael S. Bernstein. 2023. {``Generative
Agents: Interactive Simulacra of Human Behavior.''} \emph{Proceedings of
the 36th Annual ACM Symposium on User Interface Software and Technology
(UIST)}. \url{https://doi.org/10.1145/3586183.3606763}.

\bibitem[\citeproctext]{ref-perez2022discovering}
{Perez, Ethan, Sam Ringer, Kamilė Lukošiūtė, Karina Nguyen, Edwin Chen,
et al.} 2023. {``Discovering Language Model Behaviors with Model-Written
Evaluations.''} \emph{Findings of the Association for Computational
Linguistics (ACL)}.

\bibitem[\citeproctext]{ref-rose2002udl}
Rose, David H., and Anne Meyer. 2002. \emph{Teaching Every Student in
the Digital Age: Universal Design for Learning}. Association for
Supervision; Curriculum Development (ASCD).

\bibitem[\citeproctext]{ref-sharma2023sycophancy}
{Sharma, Mrinank, Meg Tong, Tomasz Korbak, et al.} 2023. {``Towards
Understanding Sycophancy in Language Models.''} \emph{arXiv Preprint
arXiv:2310.13548}.

\bibitem[\citeproctext]{ref-wiggins2005understanding}
Wiggins, Grant, and Jay McTighe. 2005. \emph{Understanding by Design}.
2nd ed. Association for Supervision; Curriculum Development (ASCD).

\bibitem[\citeproctext]{ref-wu2023autogen}
Wu, Qingyun, Gagan Bansal, Jieyu Zhang, et al. 2023. {``AutoGen:
Enabling Next-Gen LLM Applications via Multi-Agent Conversation.''}
\emph{arXiv Preprint arXiv:2308.08155}.

\end{CSLReferences}




\end{document}
